\subsection{紧空间中的连通集}

定义 3 设 V 是一致空间 X 的对称一致邻域。X 的点的有限序列 \((x_{i})_{0\leqslant i\leqslant n}\) 称为一条 V-链，如果 \(x_{i}\) 与 \(x_{i+1}\) 对 \(0\leqslant i<n\) 是 V-邻近的。点 \(x_{0}\) 与 \(x_{n}\) 称为该 V-链的端点，并称它们由这条 V-链连接。

给定对称一致邻域 V，关系“存在一条连接 x 与 y 的 V-链”是 X 中 x 与 y 之间的一个等价关系。设 \(A_{x,v}\) 是 x 关于这个关系的等价类，即可由一条 V-链与 x 连接的所有 \(y \in X\) 组成的集。显然，若 \(y \in A_{x,v}\)，则 \(\mathrm{V}(y) \subset \mathrm{A}_{x,\mathrm{v}}\)；因此 \(A_{x,v}\) 在 X 中是开的；而 \(A_{x,v}\) 的余集作为一些等价类的并，也是开的。于是：

命题 5 在一致空间 X 中，可与给定点 x 由 V-链连接的点组成的集 \(A_{x,v}\) 在 X 中既开又闭。

对每个 \(x \in X\)，设 \(A_{x}\) 表示当 V 取遍 X 的对称一致邻域时诸集合 \(A_{x,v}\) 的交；\(A_{x}\) 是 x 关于等价关系“对一切对称一致邻域 V，存在一条连接 x 与 y 的 V-链”的等价类。

命题 6 在紧空间 X 中，x 的连通分支、集合 \(A_{x}\)、以及 x 的既开又闭的邻域的交，这三者是同一个集合。

只需证明 \(A_{x}\) 是连通的：因为在任意一致空间 X 中，x 的连通分支包含在 x 的既开又闭的邻域的交中，而这个交由命题 5 包含在 \(A_{x}\) 中。

假设 \(A_{x}\) 不连通。由于 \(A_{x}\) 是闭的，存在两个不相交的非空闭集 B 与 C，使 \(B \cup C = A_{x}\)。由 §3 的命题 4，存在 X 的一致邻域 U 使 \(U(B) \cap U(C)\) 为空。

设 W 是开的一致邻域，使 \(\hat{W} \subset U\)，并设 H 是闭集，即 \(W(B) \cup W(C)\) 在 X 中的余集。例如设 \(x \in B\)，并设 y 是 C 的一点。则对每个对称一致邻域 \(V \subset W\)，由 W 的取法，对 i 用归纳法立即看出，在 X 中连接 x 与 y 的每条 V-链 \((x_{i})_{0 \leqslant i \leqslant n}\) 必有一点在 H 中。由于按假设对每个对称一致邻域 V，x 与 y 都可由一条 V-链连接，可见 \(H \cap A_{x,v}\) 当 \(V \subset W\) 时不空。另一方面，若 \(V' \subset V\)，则显然 \(A_{x,v'} \subset A_{x,v}\)；由此可得，当 V 取遍 X 的对称一致邻域时，诸集合 \(H \cap A_{x,v}\) 在紧空间 H 中构成一个由闭集组成的滤子基。因此所有这些集合有一公共点，从而 H 与 \(A_{x}\) 相交；但这与 H 的定义矛盾。

推论 设 X 是局部紧空间，K 是 X 的紧连通分支。则 K 的既开又闭的邻域构成 K 的一个基本邻域系。

设 V 是 K 在 X 中的相对紧开邻域（第一章 § 9 第 7 目命题 10），并设 F 是它的边界。设 U ⊂ V 是一个关于 V 既开又闭的集合。则 U 在 X 中是闭的，且若此外 U 不与 F 相交，则 U 在 X 中是开的（因为这时 U ⊂ V 且 U 在 V 中是开的）。因此只需证明存在 V 的一个子集，它包含 K、不与 F 相交且关于 V 既开又闭。

假设情形不是这样：则 F 与 \(\overline{V}\) 的那些包含 K 且在 \(\overline{V}\) 中既开又闭的子集的交，在 F 中构成一个由闭集组成的滤子基；F 是紧的，故这些集合有一公共点 \(y \in F\)。但这是荒谬的；由于 \(\overline{V}\) 是紧的，K 是 \(\overline{V}\) 的一个连通分支，且由命题 6，\(\overline{V}\) 中那些包含 K 且既开又闭的集合的交恰好是 K。推论至此得证。

命题 7 设 X 是紧空间，R 是 X 上的等价关系，其类是 X 的连通分支。则商空间 X/R 是紧的且全不连通。

由第一章 § II 第 5 目命题 9 我们知道 X/R 是全不连通的；因此要证明 X/R 是 Hausdorff 的（第一章 § 10 第 4 目命题 8）。设 A 与 B 是 X 的两个不同的连通分支。由命题 6，存在 X 的对称一致邻域 U，使 A 的任何点都不能由一条 U-链与 B 的任何点连接。X 中可由一条 U-链与一点 \(x \in A\)（相应地 \(y \in B\)）连接的点组成的集 V（相应地 W）在 X 中既开又闭（命题 5）且包含 A（相应地 B）；因此这两个集合分别是 A 与 B 的开邻域，关于 R 是饱和的，且不相交。证明完毕。

